\subsection{Bối cảnh và lý do chọn đề tài}

Đề tài hướng đến bài toán hỗ trợ giao tiếp cho cộng đồng người khiếm thính thông qua việc nhận diện ngôn ngữ ký hiệu, khôi phục câu tiếng Việt hoàn chỉnh và tiếp tục dịch sang tiếng Lào. Đây là một bài toán có ý nghĩa xã hội rõ ràng, đồng thời là một minh họa tiêu biểu cho cách kết hợp thị giác máy tính, xử lý ngôn ngữ tự nhiên và dịch máy trong một pipeline học sâu thống nhất.

Trong giao tiếp hằng ngày, rào cản lớn nhất giữa người dùng ngôn ngữ ký hiệu và người không biết ký hiệu nằm ở việc thông tin thị giác không được chuyển đổi tức thời sang văn bản hoặc tiếng nói. Một hệ thống chỉ nhận diện từng cử chỉ đơn lẻ vẫn chưa đủ để giải quyết vấn đề này, vì đầu ra dạng nhãn rời rạc thường thiếu ngữ pháp, thiếu ngữ cảnh và khó sử dụng trực tiếp trong hội thoại. Do đó, đề tài lựa chọn hướng tiếp cận nhiều tầng: nhận diện ký hiệu ở tầng thị giác, chuẩn hóa chuỗi ký hiệu thành câu tiếng Việt ở tầng ngôn ngữ, sau đó dịch sang tiếng Lào để mở rộng khả năng giao tiếp liên ngôn ngữ.

Lý do chọn đề tài xuất phát từ ba khía cạnh chính. Thứ nhất, đề tài có tính nhân văn vì hướng đến hỗ trợ một nhóm người dùng có nhu cầu giao tiếp đặc thù. Thứ hai, đề tài có tính ứng dụng vì có thể triển khai thử nghiệm với webcam và các mô hình học sâu hiện có. Thứ ba, đề tài có tính liên ngành cao: mỗi tầng xử lý đều là một bài toán học máy riêng, nhưng giá trị cuối cùng chỉ xuất hiện khi các tầng được ghép nối thành một hệ thống vận hành được từ đầu vào video đến đầu ra văn bản.

